% Variante : décalage de reliure pour l'impression.
% binding=12mm décale les marges intérieures (ocots-book en twoside
% uniquement — cf. tex/carrier/ocots-carrier-book.sty). Même corps que les
% autres variantes ; seule la ligne \usepackage change.
\documentclass[11pt,twoside]{ocots-book}
\usepackage[lang=fr, theme=ocots, binding=12mm, solutions=inline, math=analysis]{ocots}
\lstset{language=[LaTeX]TeX}

\title{Variante : reliure}
\author{Prénom \textsc{Nom}}
\date{\today}

\newcommand{\ocotsvariantnote}[1]{\begin{quote}\itshape Variante : #1\end{quote}}

\begin{document}

\begin{lstlisting}
\documentclass[11pt,twoside]{ocots-book}
\usepackage[lang=fr, theme=ocots, binding=12mm, ...]{ocots}
\end{lstlisting}
\ocotsvariantnote{\texttt{binding=12mm} (le défaut est \texttt{binding=0mm}) :
décalage de reliure — la marge intérieure augmente sur les pages impaires et
diminue sur les paires, pour compenser ce que la reliure masque à
l'impression. Sans effet hors \texttt{ocots-book} en \texttt{twoside}.}

\ocotscollectsolutions
% Corps partagé par toutes les variantes : aucune option, aucun préambule.
%
% C'est la démonstration de l'étape 2 : le même contenu, sans une ligne de
% différence, se compose en français ou en anglais, avec ou sans corrigés, en
% couleur ou en noir et blanc. Seule la ligne \usepackage change.

\chapter{Chapitre de démonstration}
\begin{chapterintro}
    Une introduction de chapitre, avec \ie un exemple de macro localisée et
    \enquote{un texte cité}.
\end{chapterintro}
\section{Une section}

\begin{theorem}{Un théorème}{demo}
    Énoncé.
\end{theorem}

\begin{proof}
    Démonstration.
\end{proof}

\begin{definition}{Une définition}{}
    Énoncé.
\end{definition}

\begin{remark}
    Une remarque.
\end{remark}

\begin{assumption}
    Une hypothèse, étiquetée automatiquement.
\end{assumption}

\begin{exercise}[points=4]
    Énoncé de l'exercice.
\solution
    Corrigé.
\end{exercise}

\ocotsprintsolutions
\end{document}
